\section{Extended Related Work}\label{app:extended-related-work}

\subsection{From RLHF to RLVR}

Modern language-model post-training inherits its optimization machinery from
policy-gradient reinforcement learning, from REINFORCE \cite{williams1992reinforce} to
trust-region and clipped-surrogate methods such as TRPO and PPO
\cite{schulman2015trpo,schulman2017ppo}. Advantage estimation, especially GAE, became
the standard variance-reduction device in continuous-control RL by learning value
functions over states or prefixes \cite{schulman2015gae}. In language modeling, these
ideas entered mainstream use through RLHF systems that optimize sequence-level rewards
derived from human preferences or reward models
\cite{christiano2017preferences,stiennon2020summarize,ouyang2022instruct}.

The recent reasoning-model wave shifted part of the field from preference-based RLHF
toward reinforcement learning with verifiable rewards (RLVR), where supervision is often
a deterministic correctness signal on math or code tasks. DeepSeekMath introduced GRPO
as a practical critic-free alternative for these settings \cite{shao2024deepseekmath},
and DeepSeek-R1 popularized large-scale RL-only or RL-dominant reasoning pipelines
\cite{deepseekai2025deepseekr1}. Subsequent surveys document how quickly this regime
became the default template for open reasoning-model replication and extension
\cite{zhang2025hundreddays}. This transition matters for our setting because RLVR makes
sparse outcome rewards and long reasoning trajectories central, thereby exposing the
credit-assignment problem more directly than earlier short-form RLHF tasks.

\subsection{Critic-Based and Critic-Free Policy Optimization}

The classical solution to delayed reward is to learn a critic or value function and
convert returns into token- or prefix-level advantages. PPO-style RLHF pipelines follow
this recipe, but the value-estimation problem is particularly awkward for text
generation because prefixes are high-dimensional, non-Markov, and semantically
heterogeneous. VinePPO demonstrated that standard value heads produce poor estimates of
expected returns for reasoning tasks and proposed Monte Carlo vine-style rollout
estimates as a more accurate alternative \cite{kazemnejad2024vineppo}. GenAC extends
this idea by replacing scalar value prediction with a generative critic that uses
chain-of-thought reasoning for value estimation \cite{shan2026genac}.

A growing body of work questions whether value heads are necessary at all in LLM
post-training. ReMax replaces learned critics with a greedy baseline and shows that much
of PPO's practical benefit can be retained with a simpler REINFORCE-style objective
\cite{li2023remax}. Back to Basics systematically compares PPO with RLOO-style
estimators and finds that critic-free methods can match or exceed value-based baselines
in RLHF \cite{ahmadian2024backtobasics}. REINFORCE++ further strengthens this line by
replacing the prompt-level advantage normalization used by GRPO and RLOO with a
global, batch-level normalization, stabilizing critic-free policy optimization
without reintroducing a learned value function \cite{hu2025reinforcepp}. In reasoning-focused RLVR, GRPO
normalizes rewards across a group of sampled responses rather than learning prefix
values \cite{shao2024deepseekmath}, and later work refines this family with theoretical
analyses, off-policy variants, and sequence-level formulations such as GSPO
\cite{mroueh2025revisiting,zheng2025gspo}. These methods substantially simplify the
optimization stack, but they generally still broadcast a trajectory-level signal across
all tokens once a scalar baseline is fixed.

\subsection{Reasoning-Specific Analyses of RLVR}

Another recent line asks why critic-free RLVR works so well for reasoning in the first
place. A Minimalist Approach to LLM Reasoning shows that a simple rejection-sampling
baseline (RAFT) is surprisingly competitive with GRPO and PPO on reasoning
tasks, and that GRPO's advantage over vanilla REINFORCE comes primarily from
discarding prompts whose sampled responses are all incorrect rather than from
reward normalization; the authors distill this insight into Reinforce-Rej, a
minimal REINFORCE variant that filters both entirely-correct and
entirely-incorrect samples, suggesting that much of the field's complexity is
optional rather than essential \cite{xiong2025minimalist}. Complementarily, Wen et al.\
analyze RLVR theoretically and empirically, arguing that answer-level verifiable rewards
can nonetheless incentivize correct intermediate reasoning early in a trajectory
\cite{wen2025rlvrimplicit}. These results are important because they show that uniform
outcome-level objectives already contain nontrivial reasoning signal.

At the same time, they do not eliminate the question of \emph{which} tokens should
absorb that signal. Recent empirical analysis of RLVR training dynamics suggests that
improvements are disproportionately driven by a minority of high-entropy ``forking''
tokens, and that restricting updates to this subset can remain competitive or even
outperform full-token updates in some settings \cite{wang2025beyond8020}. Building on
this observation, EAPO adapts conditional mutual information to the autoregressive RLVR
setting and proves that the credit a token can carry is upper-bounded by its entropy,
motivating an entropy-aware modulation of per-token learning signals
\cite{he2026eapo}. ERPO similarly identifies ``critical decision pivots'' at
high-entropy states and applies targeted exploration at those positions
\cite{yu2026erpo}. \emph{Sparse but Critical} argues, via a distributional-shift analysis based on
cross-sampling between policy and reference, that a small fraction of tokens
accounts for most of the useful update magnitude in RLVR, and proposes a
divergence-based selection rule for identifying them~\cite{meng2026sparsebutcritical}. These entropy- and divergence-based
methods are the closest concurrent work to the intrinsic weighting schemes that form
our baselines; the present paper does \emph{not} claim priority over them. Instead,
our contribution is to show that their underlying signals are structurally degraded
under LoRA, and to provide a credit-assignment mechanism (ARCA) that does not suffer
the same degradation.

\subsection{Fine-Grained Credit Assignment Beyond Uniform Token Updates}

A large parallel literature attempts to make supervision denser than a single
end-of-trajectory reward. One family learns explicit token- or process-level reward
estimators. Preference-grounded Token-level Guidance derives token-level signals from
preference data for fine-tuning \cite{yang2023pgtg}, while Discriminative Policy
Optimization learns token-level Q-style reward models from preferences without
requiring fine-grained annotation \cite{chen2025dpoqrm}. PRIME pushes this direction
further for reasoning tasks by constructing implicit process rewards and updating
process reward models online from rollouts and outcome labels \cite{cui2025prime}.
Reinforced Token Optimization (RTO) frames RLHF as a token-level MDP and uses DPO
log-likelihood ratios as a per-token reward signal derived from the preference
model, which is then optimized with PPO to produce fine-grained token-level
credit \cite{zhong2025rto}. These methods can produce rich
token-level feedback, but they rely on an auxiliary reward-modeling component that our
approach intentionally avoids.

A second family learns a small token-level critic on top of the policy's own hidden
states (rather than an external reward model) and uses its predictions to form
token-level advantages. This corresponds closely to an earlier version of our own
framework, and we found through an extensive literature review that the token-level
actor--critic design space was already well-covered: VinePPO~\cite{kazemnejad2024vineppo}
and GenAC~\cite{shan2026genac} cover non-parametric and generative critics
respectively, and OTB~\cite{li2026otb} provides a principled variance-minimizing
position-dependent baseline. We therefore do \emph{not} position a critic-based
configuration as the proposed method, and the submission-time sweep focuses on
critic-free intrinsic weighting schemes.

A third family redistributes outcome rewards algorithmically rather than by training a
dense reward model. RED derives token-level rewards by redistributing holistic
reward-model scores over a trajectory \cite{li2024red}. SCAR uses Shapley-inspired
attribution to estimate token or span contributions from sequence-level feedback
\cite{cao2025scar,shapley1953value}. TEMPO uses groups of sampled solutions to build a
prefix tree and compute nonparametric prefix values, yielding branch-sensitive
corrections without a learned critic \cite{tran2025tempo}. GRPO-$\lambda$ introduces
eligibility traces and lambda-returns into critic-free RLVR to propagate outcome
information backward along the sequence \cite{parthasarathi2025grpolambda}. OTB derives
an optimal position-dependent baseline that minimizes per-token gradient
variance; the practical estimator sets the baseline at position $t$ to a
cumulative-gradient-energy-weighted average of group returns-to-go, where the
weights are a causal logit-gradient proxy computed directly from the
forward-pass probabilities~\cite{li2026otb}. PSPO applies potential-based reward
shaping theory to construct dense token rewards from outcome signals without altering
the optimal policy \cite{hu2026pspo}. Compared with these approaches, our focus is
orthogonal: we do not propose a new redistribution rule, but instead diagnose a failure
mode that can affect intrinsic-signal-based redistribution methods when
applied inside LoRA.

\subsection{Alternative Action Granularities and Token Selection}

Several papers attack the same underlying problem by changing the optimization unit
rather than by redesigning the reward. MA-RLHF introduces macro actions, grouping tokens
into higher-level units to shorten the temporal distance between decisions and rewards
\cite{chai2025marlhf}. GSPO moves importance weighting and clipping from the token level
to the sequence level for greater stability in large-scale RL training
\cite{zheng2025gspo}. These methods reinforce the broader point that the granularity of
optimization matters as much as the reward definition.

Outside on-policy RL proper, token-selective optimization has also been explored in
adjacent paradigms. ConfPO emphasizes preference-critical tokens based on policy
confidence in preference optimization \cite{yoon2025confpo}, and token weighting for
long-range language modeling shows that non-uniform token emphasis can improve
optimization even in supervised settings \cite{helm2025tokenweighting}. Together with
high-entropy token analyses in RLVR \cite{wang2025beyond8020}, these results strongly
suggest that uniform token treatment is an arbitrary design choice rather than a
necessity.

\section{Additional Method Context}\label{app:method-context}

\subsection{Relationship to Existing Methods}

Our formulation recovers standard critic-free RL as a special case: uniform
weighting plus a prompt-level batch baseline yields the usual RLOO/GRPO-style
estimator. The intrinsic weighting schemes (surprisal, entropy reduction,
policy divergence) are the natural baselines within our framework and map
cleanly to published methods (e.g., divergence weighting corresponds to the
divergence-advantage analyses of \emph{Sparse but
Critical}~\cite{meng2026sparsebutcritical}, entropy-based weighting
to~\cite{he2026eapo,yu2026erpo}). ARCA is the only scheme we
study that uses an explicitly \emph{adaptation-aware} signal. Unlike
critic-based PPO, ARCA does not learn a value head. Unlike RED, PRIME, or
token-level reward-model methods, ARCA does not construct dense rewards from
an auxiliary model. Unlike TEMPO or related tree-based methods, ARCA does not
build prefix graphs or estimate branch values. Unlike hard token-selection
methods based on top-entropy masking, ARCA retains a dense token update but
modulates it continuously using an intrinsic salience score derived from the
adapter itself.

The resulting estimator is intentionally minimal. It changes only the
within-trajectory allocation of a scalar on-policy advantage, using
quantities already available from the forward pass with the adapter disabled.
This makes it easy to layer on top of existing critic-free RLVR pipelines
while preserving their computational simplicity.
